% ============================================================================
% Threats to Validity / Hạn chế — dán vào chương Đánh giá hoặc Kết luận.
% Soạn 2026-06-24 từ docs/EVAL_FINDINGS.md. Số liệu đã đo lại live trên cấu hình
% hiện tại (backbone MuQ, lyrics e5-large, nhãn emotion_labels_v6i).
% ============================================================================

\section{Hạn chế và các mối đe doạ đến tính hợp lệ (Threats to Validity)}
\label{sec:threats}

Hệ thống được đánh giá hoàn toàn bằng các chỉ số ngoại tuyến (offline). Theo
Dacrema và cộng sự (2021), tương quan giữa chỉ số offline và mức độ hài lòng thực
tế của người dùng chỉ vào khoảng $r \approx 0.28$; do đó các kết quả dưới đây là
\emph{đại lượng thay thế} (proxy), không phải bảo chứng trải nghiệm người dùng.
Chúng tôi nêu rõ các hạn chế sau để bảo đảm tính trung thực khoa học.

\subsection{Đối với gợi ý theo bài hát (similar-song)}
\begin{itemize}
  \item \textbf{Ground truth dạng playlist biên tập là thước đo không phù hợp.}
        Trên tập GT playlist biên tập, một baseline \emph{popularity} (10 bài theo
        tần suất nghệ sĩ) đạt NDCG@10 $=0.091$, cao hơn hệ sản phẩm ($0.066$). Đây
        không phải do hệ kém: GT loại này thưởng cho độ phổ biến và đồng-xuất-hiện,
        đồng thời \emph{phạt} chính cơ chế đa dạng nghệ sĩ (MMR) và lọc bản cover mà
        hệ cố ý áp dụng (Dacrema 2019). Vì vậy chúng tôi dùng GT \emph{tương tự âm
        nhạc} (được chấm theo thang 0--3) cùng các chỉ số nội tại (MoodCoherence
        $=0.94$, TempoCoherence $=0.84$) làm thước đo chính, thay vì NDCG biên tập.
  \item \textbf{Độ tin cậy của ground truth tương tự âm nhạc.} GT được kiểm chứng
        bằng \emph{hai bộ chấm độc lập} (qwen3:8b cục bộ và gpt-4o-mini của OpenAI):
        trên 1048 cặp, hai bộ chấm đồng thuận gần như tuyệt đối (Spearman
        $\rho = 0.946$, Cohen $\kappa = 0.958$) --- GT không phải tạo tác của một bộ
        chấm. Hạn chế còn lại là cỡ mẫu ($n=52$ hạt giống); tăng số hạt giống để có
        thêm công suất thống kê là hướng phát triển.
  \item \textbf{Lựa chọn backbone âm thanh.} So sánh công bằng (chuyển backbone
        thật + tối ưu lại trọng số fusion riêng cho mỗi backbone bằng 5-fold CV trên
        GT tương tự âm nhạc) cho thấy MuQ ngang-hoặc-hơn MERT: hoà trên NDCG CV
        ($0.7345$ vs $0.7343$) và audio-only ($0.7476$ vs $0.7424$, không ý nghĩa);
        MuQ \emph{vượt trội có ý nghĩa thống kê} ở phép chuyển giao sang bộ chấm độc
        lập (GPT-GT: $0.806$ vs $0.771$, CI $[+0.014, +0.056]$); color-TE coi như
        hoà. Lựa chọn MuQ là hợp lý. Lưu ý phương pháp: trọng số CV-tối-ưu hội tụ về
        gần như chỉ-âm-thanh (lyrics/V-A đóng góp $\approx 0$ vào NDCG trên pool đã
        chấm); trọng số sản phẩm $0.76/0.16/0.08$ đánh đổi một chút NDCG để tăng độ
        nhất quán cảm xúc (MoodCoherence).
\end{itemize}

\subsection{Đối với gợi ý theo màu (recommend-by-color)}
\begin{itemize}
  \item \textbf{Chỉ số Targeting Error mang tính tự quy chiếu một phần.} Nhãn V-A
        vừa là đầu vào của bộ gợi ý vừa là chuẩn so trong cùng không gian phân vị,
        nên TE thấp ($\approx 0.0242$) chỉ là điều kiện cần. Bằng chứng \emph{độc lập}
        mạnh hơn là: tương quan giữa độ sáng/độ bão hoà của màu với BPM thực đo
        (librosa) đạt $\rho = +0.41$ và $+0.66$ (Whiteford 2018), và tương quan
        valence của màu với chuẩn người ICEAS đạt $r = 0.97$.
  \item \textbf{Cơ sở dữ liệu màu--cảm xúc không bao gồm Việt Nam.} Chuẩn ICEAS
        (Jonauskaite 2020) khảo sát 30 quốc gia nhưng không có Việt Nam, và chỉ neo
        ở $n=12$ màu cố định, dẫn tới khoảng tin cậy rộng (valence $r$ Fisher-z
        $[0.89, 0.99]$; arousal chỉ $0.69$). Không tồn tại bộ dữ liệu màu--cảm xúc
        bản địa tiếng Việt để hiệu chỉnh.
  \item \textbf{Một số màu cho cảm giác phản trực giác.} Màu nâu/hồng ánh xạ sang
        arousal trung-cao theo mô hình Whiteford (ấm + bão hoà $\leftrightarrow$
        nhanh); đây là sự trung thành với mô hình màu$\leftrightarrow$nhạc, không
        phải lỗi.
  \item \textbf{Một số góc phần tư cảm xúc thưa dữ liệu.} Tương quan dương giữa
        valence và arousal trong tập nhạc khiến các vùng ``bình yên'' (V cao, A thấp)
        và ``căng thẳng'' (V thấp, A cao) ít bài, nên màu ấm phải nhắm vào phần đuôi
        arousal cao mỏng ($\approx 22\%$ tập nhạc). Đây là thuộc tính của dữ liệu.
\end{itemize}

\subsection{Hạn chế chung}
\begin{itemize}
  \item \textbf{Chưa có nghiên cứu nghe thực tế (human listening study)} do ràng
        buộc không thu thập dữ liệu người dùng; mọi chỉ số đều là proxy ngoại tuyến.
  \item \textbf{Valence chuyển giao yếu qua biên giới văn hoá.} Trong phạm vi
        nhạc phương Tây, phép chuyển giao DEAM$\to$PMEmo khá tốt ($\rho = 0.69$ cho
        valence, $0.65$ cho arousal --- đo lại 2026-06-24). Tuy nhiên khi áp bộ dò
        audio huấn luyện trên DEAM sang tập nhạc Việt thì valence rất yếu
        ($R^2 \approx 0.06$), nên valence chủ yếu dựa vào từ điển cảm xúc tiếng Việt
        ($84\%$ trọng số) thay vì âm thanh --- một lựa chọn đúng nhưng phụ thuộc
        chất lượng từ điển.
\end{itemize}
